Um zu zeigen, dass der infinitesimale Shift $\epsilon^{\alpha \beta}$ antisymmetrisch ist, betrachten wir die Erhaltung der Norm eines Vektors unter einer infinitesimalen Transformation. Sei $x^\alpha$ ein Vektor im Raum, dessen Norm durch $x^\alpha x_\alpha$ gegeben ist. Eine infinitesimale Transformation des Vektors kann durch $x'^\alpha = x^\alpha + \epsilon^{\alpha \beta} x_\beta$ beschrieben werden, wobei $\epsilon^{\alpha \beta}$ klein ist.

Die Norm des transformierten Vektors ist dann:

\[
x'^\alpha x'_\alpha = (x^\alpha + \epsilon^{\alpha \beta} x_\beta)(x_\alpha + \epsilon_{\alpha \gamma} x^\gamma).
\]

Entwickeln wir dies bis zur ersten Ordnung in $\epsilon^{\alpha \beta}$, erhalten wir:

\[
x'^\alpha x'_\alpha = x^\alpha x_\alpha + 2 \epsilon^{\alpha \beta} x_\beta x_\alpha.
\]

Da die Norm erhalten bleibt, muss gelten:

\[
x'^\alpha x'_\alpha = x^\alpha x_\alpha.
\]

Daraus folgt:

\[
2 \epsilon^{\alpha \beta} x_\beta x_\alpha = 0.
\]

Da dies für beliebige Vektoren $x^\alpha$ gelten muss, folgt, dass der symmetrische Teil von $\epsilon^{\alpha \beta}$ verschwinden muss. Das bedeutet:

\[
\epsilon^{\alpha \beta} = -\epsilon^{\beta \alpha}.
\]

Somit ist $\epsilon^{\alpha \beta}$ antisymmetrisch.